% !TEX TS-program = xelatex
% !TEX encoding = UTF-8 Unicode
% !Mode:: "TeX:UTF-8"
% Tailored for: Ecological / Environmental Data Analyst roles

\documentclass{my-resume}
\usepackage{my-resume-zh}
\usepackage{linespacing_fix}
\usepackage{cite}

\begin{document}
\pagenumbering{gobble}

\photoheader{%
  \name{吴筱琦}
  \contactInfo{+86 13776020628}{yukiwxq@gmail.com}{现居城市：伦敦/苏州}{意向城市：上海/苏州/杭州 \textbar\ 生态与环境数据分析}
}{../template/photo.jpg}{2.4cm}

\section{教育背景}
\entry{帝国理工学院 (QS 2)}{2025.09 - 2026.09}{计算生命科学研究型硕士}{伦敦}
\begin{itemize}[parsep=0.2ex]
  \item \textbf{课程内容}：生态学与进化生物学中的定量方法，包括数据科学、\textbf{GIS}、基因组学与生物信息学
\end{itemize}
\entry{帝国理工学院 (QS 2)}{2022.09 - 2025.07}{生物科学 本科}{伦敦}
\begin{itemize}[parsep=0.2ex]
  \item \textbf{相关课程}：生态学与进化论，生态学基础，行为生态学，生态学实地技能，全球变化生物学，生物多样性与保护生物学，编程与统计
\end{itemize}

\section{实习与社会实践}
\entry{上海宣泰医药科技有限公司}{2024.07 - 2024.09}{分析工程师 \textperiodcentered\ 分析部}{上海}
\begin{itemize}[parsep=0.2ex]
  \item \textbf{数据分析}：独立完成样本前处理、HPLC 系统操作及数据分析，在严格的质控流程下确保检测结果的准确性和可靠性，积累了跨领域量化数据分析经验
\end{itemize}

\section{科研与项目经历}

\entry{多模态大语言模型提取媒介性状数据评估：高精确率与结构依赖型记录召回}{2025.12 - 2026.08}{}{伦敦}
\role{Python、Gemma 4 31B IT（多模态LLM）、PyMuPDF、Docling、提示工程、JSON Schema 校验、pytest}{}
\begin{itemize}[parsep=0.2ex]
  \item \textbf{基准构建与提取流程}：设计加权多样性贪心算法从 VecTraits 超4万行、334篇论文中选取20篇，构建2,335条开发基准与431条留出测试集；结合 PyMuPDF/Docling 证据提取驱动 Gemma 4 31B IT 完成性状记录的多模态提取，为生态学文献数据整理提供可复用的量化质控框架
  \item \textbf{严谨评估}：设计二分图记录匹配评估框架，在留出测试集上验证锁定流程，取得记录精确率 1.000、条件宏观字段 F1 0.729（开发集经11轮提示词迭代达 0.810）
\end{itemize}

\entry{气候变化对佛罗里达州埃及伊蚊种群动态的影响建模}{2025.03 - 2025.06}{}{伦敦}
\role{R、deSolve、ggplot2、ODEs、统计验证（RMSE、Pearson、Spearman）}{}
\begin{itemize}[parsep=0.2ex]
  \item \textbf{模型适配与验证}：将已发表的热带 ODE 房室模型适配用于模拟佛罗里达 Collier 县埃及伊蚊种群在温度/降雨影响下的动态变化；基于 53 周监测数据验证模型（RMSE=1.65，Pearson r=0.29，Spearman $\rho$=0.30）
  \item \textbf{研究成果}：优化时间滞后（$\tau=-2$ 周）与降雨尺度因子，提出模型改进方向（人为水源、空间异质性、不确定性量化）
\end{itemize}

\entry{蜂鸟宏观生态学分析：纬度梯度与生态学规则验证}{2024.10 - 2024.12}{}{伦敦}
\role{R、ggplot2、caper、GLM、PGLS}{}
\begin{itemize}[parsep=0.2ex]
  \item \textbf{假设检验}：基于 AVONET 数据与系统发育树，应用 GLM 与 PGLS 检验纬度多样性梯度（LDG）、伯格曼法则和拉波波特法则；结果强力支持 LDG，且伯格曼法则在控制系统发育信号（PGLS）后由显著转为不显著，体现了区分统计相关性与系统发育混杂因素的方法论严谨性
\end{itemize}

\entry{环境花色比例对昆虫颜色选择的影响}{2024.09 - 2024.10}{}{开普敦}
\role{野外实验、R（t 检验、Pearson 相关）}{}
\begin{itemize}[parsep=0.2ex]
  \item \textbf{实验设计}：在南非野外课程中，于 6 个样地设计五色陷阱盘实验，研究昆虫访花颜色偏好与环境花色比例的关系
  \item \textbf{实验结果}：黄色、蓝色和白色陷阱均显著高于基线吸引昆虫（p$<$0.001）；黄色捕获率与环境黄花比例显著正相关（r=0.369，p=0.027）
\end{itemize}

\section{技能/其他}
\entrySkillsSep
\begin{itemize}[parsep=0.2ex]
  \item \textbf{数据分析与建模}：R（ggplot2、dplyr、caper、ape、vegan、deSolve）、Python、统计推断（GLM/泊松回归、PGLS、Wilcoxon 检验、t 检验、Pearson/Spearman 相关）、基于微分方程与差分方程的种群动力学建模
  \item \textbf{生态与野外研究}：野外实验设计、生物多样性/环境调查、基于 GIS 的空间分析、宏观生态学与系统发育比较分析（AVONET 级性状数据、PGLS）
  \item \textbf{语言}：英语（雅思 7.5），普通话（母语），法语（A2）；\textbf{兴趣爱好}：钢琴，摄影，滑雪
\end{itemize}

\end{document}
